Fix the public experiment \(\mathfrak E\). By \(\text{def:common-experiment}\), the supports, feature map, reference policy, and inverse temperature are public, whereas \(\rho\) is a coordinate of the law and is not fixed by \(\mathfrak E\). The first assertion is exactly the fixed-design part of \(\text{thm:fixed-design-feasibility-and-frontier-nonuniqueness}\): the variables in each branch of the certificate include \(\vartheta\), the positive context masses \((\rho_x)\), and a unit generalized-eigenvalue witness \(u\). Hence the certificate tests existence of a law-specific context distribution together with a bounded linear reward that attains both exact equalities \(C_P=C\) and \(D_P=D\). The branch enumeration and all constraints are displayed in \(\text{lem:fixed-experiment-shell-certificate}\), so this is a finite, explicit feasibility characterization rather than an unspecified existential optimization.

Suppose next that \(\mathcal M_{\mathfrak E}(d,C,D,\eta)\ne\varnothing\). The risk bound in \(\text{thm:localized-upper}\), with the public-experiment dependence restored, is
\[
\sup_{P\in\mathcal M_{\mathfrak E}(d,C,D,\eta)}
\mathbb E_{P^{\otimes n}}\!\left[J_{\eta,P}(\pi_P^\star)-J_{\eta,P}(\widehat\pi)\right]
\le
\widetilde O\!\left(\min\left\{\frac{\eta d^2D}{n},\sqrt{\frac{d^2D}{n}}\right\}\right)+n^{-2}.
\]
The learner \(\widehat\pi\) is measurable in the sample and \(\mathfrak E\), and its construction uses neither \(C\) nor \(D\). Since \(\mathfrak R_{n,\mathfrak E}\) is the infimum over all such learners, its value is no larger than the displayed risk. Inverting the fast and slow terms, as done in \(\text{thm:localized-upper}\), gives
\[
N^\star_{\mathfrak E}(\epsilon;d,C,D,\eta)
\le\widetilde O\!\left(d^2D\min\{\eta/\epsilon,\epsilon^{-2}\}\right)
\]
whenever the right-hand side is at least two. At \((C,D)=(1,1)\), \(\text{thm:zero-risk-boundary}\) instead gives the exact equality \(N^\star_{\mathfrak E}(\epsilon;d,1,1,\eta)=1\) whenever the boundary shell is nonempty.

We now show why the indices cannot determine a sharper fixed-design frontier. Fix any \(1<D\le C<e^\eta\). By \(\text{thm:fixed-design-feasibility-and-frontier-nonuniqueness}\), there is a finite public experiment \(\mathfrak E_0\) whose exact shell at these indices is nonempty and whose publicly computable Gibbs optimum is common to every law in the shell. The learner returning that policy has zero regret under every member, so
\[
\mathfrak R_{n,\mathfrak E_0}\bigl(\mathcal M_{\mathfrak E_0}(d,C,D,\eta)\bigr)=0
\quad\text{and}\quad
N^\star_{\mathfrak E_0}(\epsilon;d,C,D,\eta)=1
\]
for every \(n\in\mathbb N\) and every \(\epsilon>0\).

At the same indices, \(\text{prop:informative-fast-shell-all-indices}\) supplies another finite public experiment \(\mathfrak E_1\), constants \(c(C,D,\eta)>0\) and \(\epsilon_0(C,D,\eta)>0\), and the lower bound
\[
N^\star_{\mathfrak E_1}(\epsilon;d,C,D,\eta)
\ge c(C,D,\eta)\frac{dD\eta}{\epsilon},
\qquad 0<\epsilon\le\epsilon_0(C,D,\eta).
\]
This proposition assumes only \(1<D\le C<e^\eta\); therefore it also covers \(2D-1\ge e^\eta\). Choosing \(\epsilon\) smaller if necessary makes the last lower bound strictly larger than one, whereas the value under \(\mathfrak E_0\) remains one. Thus two experiments with the same \((d,C,D,\eta)\) have different minimax sample complexities, proving index insufficiency.

Finally, let \(L(\epsilon;d,C,D,\eta)\) be any putative lower bound required to hold for every finite public experiment with a nonempty exact shell at the displayed indices. Applying it to \(\mathfrak E_0\) yields
\[
L(\epsilon;d,C,D,\eta)
\le N^\star_{\mathfrak E_0}(\epsilon;d,C,D,\eta)=1
\]
for every \(\epsilon>0\). Hence such a bound cannot exceed one and cannot diverge as \(\epsilon\downarrow0\). In particular it cannot have order \(dD\epsilon^{-2}\). The only proved lower rate for the specified informative construction is the displayed local fast branch, and the available general upper bound retains the extra factor \(d\); therefore neither a within-experiment slow branch nor a sharper all-learner upper bound follows from these conclusions.